\documentclass[10pt,a4paper]{amsart}
\usepackage[margin=19mm]{geometry}
\usepackage{amsmath,amssymb,mathtools}
\usepackage[scr=rsfs,cal=euler]{mathalfa}
\usepackage{xfrac}
\usepackage[unicode,hidelinks,bookmarks=false]{hyperref}
\newcommand{\ie}{i.e.\ }
\newcommand{\cf}{cf.\ }
\newcommand{\resp}{respectively\ }
\newcommand{\Subsection}{\subsection}
\providecommand{\nobreakdash}{\nobreak\hbox{-}}
\setlength{\emergencystretch}{2em}
\multlinegap0pt
\usepackage{float}
\usepackage{amssymb}
\usepackage{xcolor}
\newtheorem{theorem}{Theorem}
\title{An average case efficient algorithm for solving two-variable linear Diophantine equations}
\author{Mayank Deora, Pinakpani Pal}
\date{Source: arXiv:2409.14052v6 (CC BY 4.0)}
\begin{document}
\maketitle

\noindent\emph{Source and licence.} An average case efficient algorithm for solving two-variable linear Diophantine equations. Version arXiv:2409.14052v6 (CC BY 4.0), distributed by arXiv under the Creative Commons Attribution 4.0 International licence, http://creativecommons.org/licenses/by/4.0/, as recorded in the arXiv licence metadata for this exact version and shown on the abstract page https://arxiv.org/abs/2409.14052v6. Full licence notice: \texttt{licenses/arxiv-2409.14052-CC-BY-4.0.txt}.\par\smallskip

\noindent\textit{Editorial scope.} Counted unit: the article's theorem theo10 (source0.tex lines 441--446). The article's complete original proof of it is delivered with it (source0.tex lines 448--482, one \texttt{proof} environment of 35 lines). No mathematics is written, repaired or re-derived: the statement and the proof are the article's own lines, in the article's own notation, and no retained line is substituted or re-flowed.  No further source-local context is needed: every cross-reference of every retained line resolves inside the two delivered spans.  Nothing was omitted from any retained span, so the package carries no omission record.  No retained line contains a citation, so the wrapper carries no bibliography and no bibliography entry.  No retained line contains a figure environment, an image-inclusion command or drawing code, so no image is delivered and the package declares no asset dependency.  Editorial formatting only, in addition: an AMS \texttt{amsart} preamble that restates the article's own declarations and macros used by the retained lines, namely the environments \texttt{theorem} on separate counters, together with the standard packages those lines need.  The retained environments print in this document as Theorem 1 in order of appearance, whereas the article prints the same items with the numbering of its own full document.  The complete native source of the article as distributed in the arXiv e-print for this exact version is bundled unchanged as \texttt{sources/165592/source0.tex}.
\begin{theorem} \label{theo10}

 Assume that $R_{a,b}(c)$ is a function from all the possible values of $c$ to the number of recursive function calls incurred by the DEA-R algorithm. The fundamental period of $R_{a,b}(c)$ will be $lcm(a_2,a_3,......,a_{k+1})$.
 
    
\end{theorem}

\begin{proof}
\iffalse    
For an input $a,b,c$, if $c \in \mathbf{c_i}$ and $c \notin \mathbf{c_j}$, for all $j<i$, then DEA-R algorithm will incur $i$ recursive calls to solve $ax+by=c$.
So, value of the function at $c$, $R_{a,b}(c)=i$. It is also clear that $R_{a,b}(c+a_{i+1})$
\fi
Assume that the Diophantine equation $ax+by=c$ has integer solutions. Consider the computation at line no. 6 of DEA-R algorithm, given by 
$$
(c-a_i) \bmod a_{i+1}
$$
It must evaluate to 0 for at least one value of $i$. 
This is equivalent to the statement that, if the Diophantine equation $ax+by=c$ is solvable in integers, then $c$ must belong to at least one of the sets, $\mathbf{c_1},\mathbf{c_2},....,\mathbf{c_k}$. All the integers in a set $\mathbf{c_i}$ are expressible as $a_i+Q_ia_{i+1}$ for any integer $Q_i$. In any integer $c\in \mathbf{c_i}$, we have to add at least $a_{i+1}$, so that $c+a_{i+1}$ is in $\mathbf{c_i}$ . If we consider all the values of $c$ for which Diophantine equation $ax+by=c$ is solvable in integers, we have to add atleast $lcm(a_2,a_3,.......,a_{k+1})$ in $c$ such that if $c\in c_i$, then $c+lcm(a_2,a_3,.......,a_{k+1})\in \mathbf{c_i}$. So, if $c \in \mathbf{c_i}$ but $c \notin \mathbf{c_j}$ for any $j<i$, then $c+lcm(a_2,a_3,.......,a_{k+1})\in \mathbf{c_i}$ and $c+lcm(a_2,a_3,.......,a_{k+1})\notin \mathbf{c_j}$ for any $j<i$. It implies that if DEA-R algorithm incurs $i$ recursive calls on input $a,b,c$, i.e. $R_{a,b}(c)=i$ then $R_{a,b}(c+lcm(a_2,a_3,.......,a_{k+1}))=i$.  Thus, $lcm(a_2,a_3,.......,a_{k+1})$ is a period of the function $R_{a,b}$. 
\par Now, we consider the unsolvable instances of $a,b,c$. If a $c$ does not belong to any sets $\mathbf{c_1}, \mathbf{c_2},.....,\mathbf{c_k}$, then $ax+by=c$ has no integer solutions. On the non-solvable instances of $a,b,c$, the DEA-R algorithm incurs $k+1$ recursive calls, i.e. $R_{a,b}(c)=k+1$. We can say that, $$R_{a,b}(c+lcm(a_2,a_3,.......,a_{k+1}))=k+1$$ Suppose if this is not correct, then $c+lcm(a_2,a_3,.......,a_{k+1})$ must belong to one of the sets, $\mathbf{c_1}, \mathbf{c_2},.....,\mathbf{c_k}$. So, we will be able to write $c+lcm(a_2,a_3,.......,a_{k+1})=a_i+Q_ia_{i+1}$ for any integer $Q_i$. Thus, we will be able to write $c$ also as the similar linear combination of integers $a_i$ and $a_{i+1}$. Consequently, $c$ must belong to the same sets, to which  $c+lcm(a_2,a_3,.......,a_{k+1})$ belong. So, this is a contradiction and for unsolvable instances of $a,b,c$ also, $R_{a,b}(c)=R_{a,b}(c+lcm(a_2,a_3,.......,a_{k+1}))$. 
\par Hence, we proved that $lcm(a_2,a_3,.......,a_{k+1})$ is a period of the function $R_{a,b}$. But, to prove that this is the fundamental period, we have to consider the case when there is a set $\mathbf{c_i}$, such that for all $c \in \mathbf{c_i}$, there is a set $\mathbf{c_j}$ such that $c \in \mathbf{c_j}$, where $j<i$. In this case, the period of $R_{a,b}$ may be smaller than $lcm(a_2,a_3,.......,a_{k+1})$, because $R_{a,b}$ will never be equal to $i$, so that we can divide the $lcm(a_2,a_3,.......,a_{k+1})$ by $a_{i+1}$ to get a smaller period. We will prove that this scenario is not possible for $i<k$ i.e. there exists atleast one value in each of the sets $\mathbf{c_1},\mathbf{c_2},..., \mathbf{c_{k-1}}$ which does not belong to any smaller indexed set. Specifically, we will prove that for all $i<k$, there exists a $c$ with $c \in \mathbf{c_i}$ but $c \notin \mathbf{c_j}$ for all $j<i$. After this we will see that for set $\mathbf{c_k}$, we do not need an integer which belongs to only the set $\mathbf{c_k}$.    
    \par  Now we will prove that, for all $i<k$, $c=a_{i+2}$ belongs to $\mathbf{c_i}$ but it does not belong to $\mathbf{c_j}$ for all $j<i$. Consider the integer $$c=a_{i+2}$$
$$\implies c = a_i-q_ia_{i+1}$$ 

where $q_i$ (defined in notations section) is the quotient of division $\frac{a_i}{a_{i+1}}$.
$$\implies c \in \mathbf{c_i}$$
because all the integers belonging to $\mathbf{c_i}$ are in the form $a_i+Q_ia_{i+1}$ (see notations Section).
Suppose that $a_{i+2} \in \mathbf{c_j}$ for some $j<i$, then 
\begin{equation}  \label{perEq1}
a_{i+2}=a_j+Q_{j}a_{j+1}    
\end{equation}




for an integer $Q_j$, which is positive, negative or $0$. 
\par $j<i$ implies that  $a_i<a_j$ and since $a_{i+2}<a_i$, so $a_{i+2}<a_j$ and $a_{i+2}<a_{j+1}$. Since $a_j$ and $a_{j+1}$ are positive integers, for Equation (\ref{perEq1}) to hold, $Q_j$ must be negative. We write equation \ref{perEq1} as follows:
\begin{equation}\label{perEq2}
    a_{i+2}-Q_ja_{j+1}=a_j
\end{equation}
Since $-Q_j$ is a positive integer and $a_{i+2}<a_{j+1}$, $a_{i+2}$ is the remainder of the division, $\frac{a_j}{a_{j+1}}$. It is a contradiction because we know that $a_{i+2}$ is the remainder of the division, $\frac{a_i}{a_{i+1}}$. Thus for all $i<k$, $a_{i+2}\in \mathbf{c_i}$ but $a_{i+2} \notin \mathbf{c_j}$ for all $j<i$. 
\par Suppose that we have a value in $\mathbf{c_k}$ which belongs to only $\mathbf{c_k}$ and no other sets $\mathbf{c_j}$ such that $j<k$. Then the fundamental period of $R_{a,b}$ will be $lcm(a_2,a_3,.....,a_{k+1})$. Now, assume that all the values in set $\mathbf{c_k}$ belong to any of the sets $c_j$ for some $j<k$. Then fundamental period of $R_{a,b}$ will be $lcm(a_2,a_3,......,a_{k})$. Since $a_{k+1}$ is the $gcd$ of $a_2$ and $a_3$, $lcm(a_2,a_3,......,a_{k+1})=lcm(a_2,a_3,.......,a_k)$. Thus the fundamental period of function $R_{a,b}$ is equal to $lcm(a_2,a_3,......,a_{k+1})$.
\end{proof}

\end{document}
